\section{基准测试（Benchmarking）}

\subsection{为什么要基准测试}

基准测试测量操作的\textbf{端到端墙钟时间}（wall-clock time）。虽然简单，但它是性能优化的基础：

\begin{itemize}
    \item 比较不同实现的速度（``CUDA 版本比 PyTorch 快吗？''）
    \item 理解性能如何随参数缩放（``矩阵增大一倍，时间增加多少？''）
\end{itemize}

\begin{warningbox}{不能只靠理论推断性能}
你可以阅读 spec sheet（营销材料）和论文，但实际性能取决于你的\textbf{库版本}、\textbf{硬件}和\textbf{工作负载}。没有什么能替代实际的基准测试。
\end{warningbox}

\subsection{正确的基准测试方法}

一个正确的 GPU 基准测试函数需要注意以下几点：

\begin{lstlisting}[caption={基准测试函数模板（Python）}]
def benchmark(description, run, num_warmups=1, num_trials=3):
    # 1. Warmup: 避免测量启动开销
    for _ in range(num_warmups):
        run()
    if torch.cuda.is_available():
        torch.cuda.synchronize()  # 等待 GPU 完成

    # 2. 计时
    times = []
    for trial in range(num_trials):
        start_time = time.time()
        run()
        if torch.cuda.is_available():
            torch.cuda.synchronize()  # 重要！
        end_time = time.time()
        times.append((end_time - start_time) * 1000)

    return mean(times)
\end{lstlisting}

\begin{importantbox}{基准测试两大要点}
\begin{enumerate}
    \item \textbf{Warmup}：首次运行时，PyTorch 会在后台编译 CUDA 代码、发送指令到 GPU 等。这些一次性开销不应计入稳态性能。
    \item \textbf{torch.cuda.synchronize()}：GPU 和 CPU 是\textbf{两个独立的计算设备}。CPU 发送 CUDA 内核到 GPU 后会立即返回（异步执行）。如果不调用 synchronize()，你测量的是 CPU 端\textbf{入队}的时间，而非 GPU 端\textbf{执行}的时间。
\end{enumerate}
\end{importantbox}

\subsection{基准测试案例：矩阵乘法}

对方阵矩阵乘法进行基准测试，可以观察到性能随维度的变化：

\begin{table}[H]
\centering
\begin{tabular}{cc}
\toprule
\textbf{矩阵维度} & \textbf{典型耗时趋势} \\
\midrule
1024 & 基线 \\
2048 & $\sim$3--4x 增长 \\
4096 & $\sim$8--12x 增长 \\
8192 & 显著增长，接近理论 $O(N^3)$ \\
16384 & 非常大，充分利用 GPU \\
\bottomrule
\end{tabular}
\caption{矩阵乘法的性能随维度缩放}
\end{table}

\subsection{基准测试案例：MLP 缩放}

用一个简单的 MLP（Linear $\rightarrow$ GeLU $\rightarrow$ Linear $\rightarrow$ GeLU $\rightarrow \ldots$）作为实验对象，分别缩放不同参数：

\begin{lstlisting}[caption={MLP 定义}]
class MLP(nn.Module):
    def __init__(self, dim, num_layers):
        super().__init__()
        self.layers = nn.ModuleList(
            [nn.Linear(dim, dim) for _ in range(num_layers)]
        )

    def forward(self, x):
        for layer in self.layers:
            x = layer(x)
            x = torch.nn.functional.gelu(x)
        return x
\end{lstlisting}

缩放实验的发现：
\begin{itemize}
    \item \textbf{步数 (num\_steps)}：线性缩放——2x 步数 $\approx$ 2x 时间
    \item \textbf{层数 (num\_layers)}：线性缩放——每层增加固定时间
    \item \textbf{批大小 (batch\_size)}：小批量时增长缓慢（GPU 未充分利用），大批量时接近线性
    \item \textbf{维度 (dim)}：增长不太可预测——涉及矩阵乘法的非均匀 CUDA 内核调度
\end{itemize}

\begin{knowledgebox}{性能不总是可预测的}
由于 CUDA 内核调度、硬件特性、缓存行为等因素，基准测试结果不总是符合简单的理论预测。这正是实际基准测试不可替代的原因。
\end{knowledgebox}

\subsection{本章小结}

基准测试是性能优化的第一步。正确的基准测试需要 warmup 和 CUDA synchronize。它告诉你``多快''，但不告诉你``为什么快/慢''——这需要 Profiling。

%% ========== Section 4: Profiling ==========

